\documentclass[11pt]{article}
\usepackage[a4paper,margin=0.78in]{geometry}
\usepackage{amsmath,amssymb,amsthm,mathtools}
\usepackage{enumitem,microtype}
\usepackage[dvipsnames]{xcolor}
\usepackage[colorlinks=true,linkcolor=MidnightBlue,citecolor=MidnightBlue,urlcolor=MidnightBlue]{hyperref}
\setlist{nosep,leftmargin=*}
\newtheorem{theorem}{Theorem}
\newtheorem{corollary}[theorem]{Corollary}
\newtheorem{proposition}[theorem]{Proposition}
\newcommand{\Hh}{\mathcal H}
\newcommand{\Pw}{\mathcal P(\Hh)}
\newcommand{\Aw}{\mathcal A_w}
\newcommand{\esssup}{\operatorname*{ess\,sup}}
\title{Hilbert-valued weighted Sobolev approximation\\
\large The missing uniform-tail condition and an exact primitive criterion}
\author{Autonomous solve-first investigation, lane 12}
\date{13 August 2026}

\begin{document}
\maketitle

\begin{abstract}
Pérez and Quintana ask for the closure of Hilbert-valued polynomials in a
weighted $W^{1,\infty}$ space.  We first repair the infinite-dimensional
$L^\infty$ input: coordinatewise scalar approximability is necessary but not
sufficient; the missing condition is uniform decay of the weighted
coordinate tails.  This gives an exact closure theorem and a counterexample
to the coordinatewise-only assertion in the companion paper.  We then prove
that whenever the weighted primitive operator is bounded, the Sobolev
closure consists exactly of the functions whose derivative satisfies that
corrected $L^\infty$ criterion.  A one-sided coordinate comparison is an
explicit sufficient condition.  In particular, for every constant positive
diagonal Hilbert weight $w$, the answer is exactly
$\{f:D_wf\in C^1(I;\Hh)\}$.  The unrestricted formulation also permits
weights for which the displayed quantity is only a seminorm, so no claim is
made for every weakly measurable weight.
\end{abstract}

\section{Corrected setting and the weighted sup-norm closure}

Let $I=[a,b]$, let $\Hh$ be a separable real Hilbert space with fixed
orthonormal basis $(e_j)_{j\ge1}$, and write $f_j=\langle f,e_j\rangle$.
Let $w=(w_j)_{j\ge1}\in L^\infty(I;\Hh)$, with $w_j>0$ almost everywhere.
Multiplication is coordinatewise, as in the source, and
\[
 \|f\|_w=\esssup_{t\in I}
 \left(\sum_{j\ge1}|f_j(t)w_j(t)|^2\right)^{1/2}.
\]
Denote the resulting space by $X_w$.  For each $j$, let
\[
 \mathcal A_j=\overline{C(I)}^{\,L^\infty(I,w_j)};
\]
the scalar weighted Weierstrass theorem quoted in \cite{PQ} describes this
space by one-sided regularity and weighted matching at the singularities of
$w_j$.  Define
\begin{equation}\label{eq:Aw}
 \Aw=\left\{f\in X_w: f_j\in\mathcal A_j\ (j\ge1),\quad
 \lim_{N\to\infty}\esssup_t
 \sum_{j>N}|f_j(t)w_j(t)|^2=0\right\}.
\end{equation}

\begin{theorem}[uniform-tail correction]\label{thm:Linf}
The closures in $X_w$ of $C(I;\Hh)\cap X_w$ and of the
$\Hh$-valued polynomials $\Pw$ both equal $\Aw$.
In finite dimension the tail condition is automatic.
\end{theorem}

\begin{proof}
Suppose first that $g\in C(I;\Hh)$.  Coordinate extraction is contractive
from $X_w$ to $L^\infty(I,w_j)$.  Moreover, compactness of $g(I)$ gives
$\sup_t\|(I-P_N)g(t)\|\to0$, where $P_N$ is coordinate projection.  Hence
\[
 \esssup_t\sum_{j>N}|g_j(t)w_j(t)|^2
 \le \|w\|_{L^\infty(I;\Hh)}^2
       \sup_t\|(I-P_N)g(t)\|^2\longrightarrow0.
\]
Both properties pass to $X_w$ limits, proving necessity.

Conversely, let $f\in\Aw$ and fix $\varepsilon>0$.  Choose $N$ so the
tail in \eqref{eq:Aw} has square root below $\varepsilon/2$.  For
$1\le j\le N$, choose $g_j\in C(I)$ such that
\[
 \|f_j-g_j\|_{L^\infty(I,w_j)}<\frac{\varepsilon}{2\sqrt N}.
\]
Then $g=\sum_{j=1}^Ng_je_j$ is continuous and $\|f-g\|_w<\varepsilon$.
Finally, vector-valued Weierstrass approximation gives $p\in\Pw$ uniformly
close to $g$, and
\[
 \|g-p\|_w\le \|w\|_{L^\infty(I;\Hh)}\|g-p\|_{L^\infty(I;\Hh)}.
\]
Thus polynomials have the same closure.
\end{proof}

The extra condition in \eqref{eq:Aw} cannot be discarded.

\begin{proposition}[coordinatewise approximation is insufficient]
Let $\Hh=\ell^2$, $I=[0,1]$, and $w_j(t)=2^{-j}$.  There is
$f\in X_w$ such that every $f_j$ is continuous, but
\[
 \operatorname{dist}_{X_w}(f,C(I;\ell^2))\ge\tfrac12.
\]
Consequently, the coordinatewise-only infinite-dimensional closure
assertion in \cite{Q} is false as written.
\end{proposition}

\begin{proof}
Put $t_j=2^{-j}$ and choose continuous triangular bumps $\phi_j$ of height
one, centered at $t_j$, with pairwise disjoint supports and with
$0$ outside every support.  Define $f_j=2^j\phi_j$.  At each $t$ at most
one coordinate is nonzero, so $f(t)\in\ell^2$, while $\|f\|_w=1$ and every
$f_j$ is continuous.  If $g\in C(I;\ell^2)$ and $\|f-g\|_w<1/2$, evaluation
at $t_j$ gives
\[
 |1-2^{-j}g_j(t_j)|<\tfrac12,\qquad
 |g_j(t_j)|>2^{j-1}.
\]
The essential-supremum inequality holds at the centers as well because both
sides of each coordinate inequality are continuous.  This contradicts
boundedness of the continuous function $g$ on $I$.
Equivalently, the uniform weighted tail of $f$ is always one.
\end{proof}

\section{Exact Sobolev lifting by the primitive operator}

To remove the representative ambiguity in the source definition, set
\[
 W_w^{1,\infty}=\{f\in AC(I;\Hh): f,f'\in X_w\},
 \qquad \|f\|_{W_w}=\|f\|_w+\|f'\|_w,
\]
where $AC$ means Bochner absolutely continuous.  Assume that the primitive
$Vg(t)=\int_a^t g(s)\,ds$ obeys
\begin{equation}\label{eq:V}
 \|Vg\|_w\le C_V\|g\|_w
\end{equation}
whenever $g\in X_w\cap L^1(I;\Hh)$.

\begin{theorem}[weighted primitive criterion]\label{thm:Sob}
Under \eqref{eq:V},
\begin{equation}\label{eq:closure}
 \overline{\Pw}^{\,W_w^{1,\infty}}
 =\{f\in W_w^{1,\infty}: f'\in\Aw\}.
\end{equation}
Thus \eqref{eq:Aw} gives an exact coordinate-and-tail characterization of
the Sobolev closure for every weight satisfying \eqref{eq:V}.
\end{theorem}

\begin{proof}
If $p_n\to f$ in $W_w^{1,\infty}$, then the polynomials $p_n'$ converge to
$f'$ in $X_w$, so Theorem~\ref{thm:Linf} gives $f'\in\Aw$.
Conversely, if $f'\in\Aw$, choose $q_n\in\Pw$ with
$\|q_n-f'\|_w\to0$, and define the polynomial
\[
 p_n(t)=f(a)+\int_a^tq_n(s)\,ds.
\]
Then $p_n'=q_n$, while absolute continuity and \eqref{eq:V} give
\[
 \|f-p_n\|_w=\|V(f'-q_n)\|_w
 \le C_V\|f'-q_n\|_w\longrightarrow0.
\]
This proves convergence in the full Sobolev norm.
\end{proof}

Here is a directly checkable weight hypothesis.  If, for some $C\ge1$,
\begin{equation}\label{eq:ratio}
 w_j(t)\le Cw_j(s)\qquad(a\le s\le t\le b, j\ge1)
\end{equation}
outside a common null set, then Minkowski's inequality gives
\[
 \|(Vg)(t)w(t)\|_{\ell^2}
 \le\int_a^t\left(\sum_j|w_j(t)g_j(s)|^2\right)^{1/2}ds
 \le C(b-a)\|g\|_w.
\]
Thus \eqref{eq:V} holds.  In particular it holds for every constant weight.

\begin{corollary}[complete constant-weight answer]\label{cor:constant}
Let $w(t)=\alpha=(\alpha_j)\in\ell^2$, where every $\alpha_j>0$, and let
$D_\alpha e_j=\alpha_je_j$.  Then
\[
 \overline{\Pw}^{\,W_w^{1,\infty}}
 =\{f\in W_w^{1,\infty}:D_\alpha f\in C^1(I;\Hh)\}.
\]
\end{corollary}

\begin{proof}
For constant $w$, Theorem~\ref{thm:Linf} says precisely that
$h\in\Aw$ iff $D_\alpha h$ has a continuous $\Hh$-valued representative.
Indeed, one direction follows by uniform limits; for the other, approximate
the finitely many coefficients of a vector polynomial by vectors in the
dense range of $D_\alpha$.  Since $D_\alpha$ is bounded and $f$ is absolutely
continuous, $(D_\alpha f)'=D_\alpha f'$ almost everywhere.  Apply
Theorem~\ref{thm:Sob}.
\end{proof}

\section{Scope and obstruction to the unrestricted question}

The source permits arbitrary weakly measurable Hilbert-valued weights.  In
that generality its displayed quantity need not be a norm: for
$\Hh=\mathbb R^2$ and $w=(1,0)$, the nonzero constant $f=(0,1)$ has
$\|f\|_w=0$.  The same weighted class contains representatives whose second
coordinate has no derivative, so Sobolev membership is not representative
independent without an additional convention.  Even after imposing
nondegeneracy, \eqref{eq:V} can fail when a coordinate weight grows too fast
relative to its earlier values.  Theorems~\ref{thm:Linf} and
\ref{thm:Sob} therefore solve a broad corrected class and isolate the two
remaining issues---well-posedness and primitive boundedness---but do not
claim the most-general classification requested in \cite{PQ}.

Eight focused upgrade routes were tested.  Exact-title, author, journal,
scalar weighted-Sobolev, and citation searches on 13 August 2026 found no
later solution of the stated Hilbert-valued Sobolev problem and no correction
of the missing uniform-tail condition.  Novelty is plausible, not certified.

\begin{thebibliography}{9}
\bibitem{PQ}
D. Pérez and Y. Quintana,
\emph{A survey on the Weierstrass approximation theorem},
Divulg. Mat. 16 (2008), 231--247; arXiv:math/0611038.

\bibitem{Q}
Y. Quintana,
\emph{On Hilbert extensions of Weierstrass' theorem with weights},
J. Funct. Spaces Appl. 8 (2010), 201--213; arXiv:math/0611034.
\end{thebibliography}

\end{document}
